\documentclass[12pt, a4paper]{article}

\usepackage[margin=1in]{geometry}
\usepackage{times}
\usepackage{parskip}
\usepackage{titlesec}
\usepackage{hyperref}
\usepackage{enumitem}
\usepackage{booktabs}
\usepackage{array}
\usepackage{xcolor}

\definecolor{sectioncolor}{RGB}{30, 80, 150}

\titleformat{\section}{\large\bfseries\color{sectioncolor}}{}{0em}{}[\titlerule]
\titleformat{\subsection}{\normalsize\bfseries}{}{0em}{}

\hypersetup{
    colorlinks=true,
    linkcolor=sectioncolor,
    urlcolor=sectioncolor
}

\title{
    \textbf{GitDB: Version Control System for MySQL Databases} \\[0.4em]
    \large Abstract and Feasibility Study \\[0.4em]
    \normalsize DISL Mini Project --- SY-11, Lab Batch F-11 \\
    Pune Institute of Computer Technology (PICT), Pune
}

\author{
    Aditya Malode (23325) \and
    Rohan Natekar (23330) \and
    Aarush Oak (23332) \and
    Prathamesh Kinagi (23340)
}

\date{Academic Year 2025--26}

\begin{document}

\maketitle
\thispagestyle{empty}

\section{Abstract}

Version control has become a foundational pillar of modern software engineering.
Systems like Git allow developers to snapshot code, inspect historical changes,
revert to any prior state, and collaborate across teams with full auditability.
Yet despite MySQL databases being a critical component of almost every production
software system, they have never enjoyed an equivalent first-class versioning
mechanism. Schema changes are typically managed through manually written migration
scripts that are fragile, difficult to audit, and entirely disconnected from the
data they govern. Data rollbacks, when needed, require either expensive full
database dumps or complex reconstruction --- neither of which resembles the
intuitive commit-and-revert workflow developers take for granted with source code.

\textbf{GitDB} addresses this gap directly. It is a MySQL-native, state-based
version control system designed to bring Git-like semantics --- \texttt{commit},
\texttt{log}, \texttt{diff}, and \texttt{checkout} --- to MySQL database management
without requiring developers to manually author SQL migration scripts. The central
idea is simple: instead of tracking \textit{instructions} for how to reach a state,
GitDB tracks the \textit{state itself}.

At every commit, a complete snapshot of the target MySQL database's schema (DDL)
and data is captured and persisted into a dedicated metadata database named
\texttt{gitdb\_meta}. This clean separation mirrors the way Git's \texttt{.git/}
directory is isolated from the source files it tracks. Each commit is identified
by a SHA-256 hash computed from its content serialisation and the hash of its
parent commit, forming an immutable, tamper-evident Merkle chain.

The system's core technical contribution is its diff engine, which takes two
commit snapshots and automatically generates the MySQL SQL necessary to transition
the live database from one state to another. Schema diffing is performed at the
Abstract Syntax Tree (AST) level using \texttt{sqlglot}, correctly identifying
added tables, dropped tables, and altered columns. Data diffing uses a
primary-key-anchored algorithm operating in $O(n)$ time.

A key engineering challenge specific to MySQL is that DDL statements issue
implicit commits and cannot be rolled back in a standard transaction. GitDB
handles this through a \textbf{two-phase checkout strategy}: schema changes are
applied and validated first; data patches are then applied inside an explicit
transaction with rollback on failure.

The system is delivered through two complementary interfaces: a command-line
interface built with \texttt{click} and \texttt{rich} providing Git-familiar
commands, and a web-based interface built in React~18 with Tailwind CSS offering
a visual commit graph and side-by-side diff viewer.

\medskip
\noindent\textbf{Keywords:} MySQL, Version Control, State-Based Snapshotting,
Schema Diffing, SQL Patch Generation, Merkle Hash Chain, DDL Auto-Commit,
Two-Phase Checkout, Python, React, Flask.

\section{Feasibility Study}

\subsection{Technical Feasibility}

GitDB is built entirely on mature, production-grade technologies with no
experimental dependencies. The following aspects confirm technical feasibility:

\begin{itemize}[leftmargin=1.5em]

    \item \textbf{MySQL Information Schema:} MySQL natively exposes full schema
    metadata via \texttt{information\_schema.TABLES}, \texttt{information\_schema.COLUMNS},
    and \texttt{SHOW CREATE TABLE}. These are stable, documented APIs that make
    schema introspection reliable without any reverse engineering.

    \item \textbf{State-based snapshotting:} Capturing schema and data as JSON
    blobs at each commit is straightforward using Python's \texttt{mysql-connector}
    and the standard \texttt{json} library. The 10,000-row per table limit keeps
    snapshot sizes manageable within a mini-project scope.

    \item \textbf{AST-level diffing:} The \texttt{sqlglot} library provides
    robust SQL parsing across dialects including MySQL, enabling structural
    comparison of DDL at the AST level rather than fragile text comparison.

    \item \textbf{Merkle hashing:} SHA-256 via Python's \texttt{hashlib} is
    deterministic and collision-resistant, making it suitable for commit
    identification and tamper detection without any cryptographic overhead.

    \item \textbf{DDL Auto-Commit workaround:} MySQL's implicit DDL commit is
    a well-documented limitation. The two-phase checkout strategy --- applying
    schema changes first and falling back to the stored snapshot on failure ---
    is a principled engineering solution that does not require patching MySQL or
    using non-standard drivers.

    \item \textbf{Web interface:} React~18, Tailwind CSS, \texttt{react-flow},
        and \texttt{@git-diff-view/react} are stable, widely adopted libraries.
    Flask provides a lightweight REST backend with minimal configuration.

\end{itemize}

\subsection{Operational Feasibility}

GitDB is designed to be operationally straightforward for its target users:

\begin{itemize}[leftmargin=1.5em]

    \item \textbf{Familiar interface:} The CLI mirrors Git's command structure
    (\texttt{init}, \texttt{commit}, \texttt{log}, \texttt{diff},
    \texttt{checkout}, \texttt{status}), reducing the learning curve for
    developers already familiar with Git.

    \item \textbf{No schema modification:} The target MySQL database is never
    structurally altered by GitDB itself. All metadata lives in the separate
    \texttt{gitdb\_meta} database, ensuring zero pollution of the versioned
    database.

    \item \textbf{Self-contained deployment:} GitDB requires only a running
    MySQL instance, Python~3.10+, and Node.js for the frontend. No additional
    infrastructure (message queues, object storage, external APIs) is needed.

    \item \textbf{Documented constraints:} Column type changes and table renames
    are explicitly flagged rather than silently mishandled, consistent with the
    design philosophy of production tools like Prisma and Atlas.

\end{itemize}


\medskip
GitDB is a well-scoped, technically grounded mini-project that demonstrates
substantive breadth across core DBMS concepts --- schema introspection, ACID
properties, DDL constraints, cryptographic integrity, and relational modelling ---
while delivering a working system that solves a real and recognised problem in
MySQL development workflows.

\end{document}